%% =====================================================================
%%  RUTA DE CIENCIA  —  Post #02
%%  Pilar: DETECTOR DE MENTIRAS  (enfocado en CRECIMIENTO / seguidores)
%%  Tema:  "Tu horóscopo SIEMPRE le atina" — efecto Barnum / Forer
%%
%%  Estrategia de crecimiento:
%%   - Hook con demostración interactiva (alta compartibilidad).
%%   - Momento "mind-blown" dentro del propio carrusel.
%%   - Lámina final de SÍGUENOS + gancho de etiquetar (alcance).
%%
%%  Compilar:  ./build.sh post-02-horoscopo-barnum
%% =====================================================================
\documentclass[10pt]{beamer}
\usepackage{ruta-de-ciencia}

% ---- Metadatos de este post ----------------------------------------
\renewcommand{\rdchandle}{@rutadeciencia}
\renewcommand{\rdccategory}{DETECTOR DE MENTIRAS}
\renewcommand{\rdctotalslides}{9}

\begin{document}

% --- 1 · PORTADA (hook) --------------------------------------------
\rdccover{MITO}
  {Tu horóscopo\\ SIEMPRE\\ ``le atina''}
  {Y no es magia: es un truco psicológico con nombre.
   Te lo demuestro en 10 segundos.}

% --- 2 · LA PRUEBA --------------------------------------------------
\rdccontent{2}{\faComments}{HAGAMOS UNA PRUEBA}
  {Lee esto con calma}
  {Voy a describir tu personalidad sin saber tu signo. En la siguiente
   lámina, cuenta cuántas frases sientes que hablan de ti.}

% --- 3 · LA "LECTURA" (texto Barnum) -------------------------------
\rdccontent{3}{\faQuoteLeft}{TU LECTURA DE HOY}
  {Esto eres tú:}
  {``Necesitas caer bien, pero eres tu crítico más duro. Tienes un
   potencial que aún no explotas. A veces dudas de tus decisiones y
   valoras tu independencia por encima de las reglas.'' \\[4pt]
   \textcolor{rdcamber}{¿Cuántas te describieron?}}

% --- 4 · EL TRUCO ---------------------------------------------------
\rdccontent{4}{\faMagic}{EL TRUCO}
  {Le atina a todos}
  {Esas frases le quedan a casi cualquier persona. Se llama efecto
   Forer (o Barnum): descripciones tan generales que sientes que hablan
   solo de ti.}

% --- 5 · POR QUÉ CAEMOS ---------------------------------------------
\rdccontent{5}{\faBrain}{POR QUÉ CAES}
  {Tu cerebro rellena}
  {Recuerdas los aciertos y olvidas los fallos (sesgo de confirmación).
   Y como quieres que sea verdad, tu mente completa los huecos por ti.}

% --- 6 · LA PRUEBA CIENTÍFICA ---------------------------------------
\rdccontent{6}{\faFlask}{¿Y SI FUERA REAL?}
  {Lo pusieron a prueba}
  {En un estudio doble ciego publicado en \emph{Nature} (Carlson, 1985),
   astrólogos no lograron emparejar cartas astrales con perfiles de
   personalidad mejor que el puro azar.}

% --- 7 · SIN DRAMA --------------------------------------------------
\rdccontent{7}{\faHeart}{SIN DRAMA}
  {Disfrútalo, pero...}
  {No hay nada malo en leer tu horóscopo por diversión. El problema es
   tomar decisiones reales --de salud, dinero o amor-- basándote en los
   astros.}

% --- 8 · VEREDICTO --------------------------------------------------
\rdcclosing{8}
  {El horóscopo no te conoce. Tú te conoces (y rellenas los huecos).}
  {La próxima vez que ``le atine'', recuerda: está diseñado para
   atinarle a todo el mundo.}
  {Forer, B.R. (1949); Carlson, S. (1985), \emph{Nature} 318, 419--425.}

% --- 9 · SÍGUENOS (crecimiento) ------------------------------------
\rdcfollow{9}
  {¿Te voló la cabeza?}
  {Cada semana desarmo un mito con evidencia. Sígueme y no te comas
   otro cuento.}
  {Y etiqueta a quien SIEMPRE saca su carta astral.}

\end{document}
